\chapter*{Introduction}
\addcontentsline{toc}{chapter}{Introduction}

In the context of a rapidly growing digital economy, the volume of financial transactions executed every day worldwide is enormous. Alongside this convenience comes an ever-increasing risk of financial fraud, particularly money laundering and the splitting of cash flows (\textit{smurfing}) used to evade monitoring thresholds. Manually inspecting every transaction is infeasible; therefore, the need to build automated software systems capable of scanning and detecting suspicious transactions has become urgent.

Such a system requires not only business-level accuracy but also a sound software architecture: easy to maintain and easy to extend as legal regulations change and new types of fraud emerge. This is an ideal setting in which to apply the principles of Object-Oriented Programming (OOP) -- a programming paradigm that helps model complex business entities as clear software objects with high reusability and extensibility.

Motivated by this reality, the group chose the topic ``Design and Implementation of a Financial Transaction Management \& Audit System using Object-Oriented Programming''. The objective is to design and implement an engine capable of managing accounts and transactions and automatically assessing the risk of each transaction based on a polymorphic set of audit rules, while fully and deeply demonstrating the object-oriented techniques studied in the course.

\newpage
The content of this report is organised into three chapters as follows:

\begin{itemize}
    \item \textbf{Chapter 1: Project Overview.} Presents the context of the financial fraud detection problem, the objectives and scope of the system, the functional requirements, and the tools and development environment used.
    \item \textbf{Chapter 2: Theoretical Background and System Design.} Introduces the fundamental principles of Object-Oriented Programming, then presents in detail the overall architecture of the system through a UML class diagram, clarifying the role of each module and how the OOP principles are applied to the design.
    \item \textbf{Chapter 3: Implementation and Evaluation.} Presents in detail the implementation of the core components, illustrated with source code, the experimental results of the system, the verification of memory management, and the division of work within the group.
\end{itemize}

\vspace{1cm}
\hspace{0.5\textwidth}
\begin{minipage}{0.5\textwidth}
	\noindent\begin{center}
		\vspace{0.5cm}
		\textit{Hanoi, \today} \\
		The group of students
	\end{center}	
\end{minipage}
